\documentclass[a4paper,10pt,twocolumn]{ltjarticle}
\usepackage{kaitoushu}

\pgfplotsset{
  myaxis/.style={
    axis lines=middle,
    xlabel={$x$}, ylabel={$y$},
    every axis x label/.style={at={(ticklabel* cs:1)},anchor=west},
    every axis y label/.style={at={(ticklabel* cs:1)},anchor=south},
    tick label style={font=\scriptsize},
    label style={font=\small},
    width=5.5cm, height=5cm,
    grid=none,
    samples=200,
  }
}

\begin{document}

\problemrange{5章：三角関数 — Check (p.57)（問262--271）}



\mondai{262}{次の問いに答えよ．}
\kotae{
$\sin^{2}\alpha+\cos^{2}\alpha=1$ で残りを求める．

(1) $\alpha$鋭角，$\sin\alpha=\dfrac{1}{4}$：$\cos\alpha=\dfrac{\sqrt{15}}{4}$，$\tan\alpha=\dfrac{1}{\sqrt{15}}=\dfrac{\sqrt{15}}{15}$

(2) $\alpha$鈍角，$\sin\alpha=\dfrac{1}{4}$：$\cos\alpha=-\dfrac{\sqrt{15}}{4}$，$\tan\alpha=-\dfrac{\sqrt{15}}{15}$

(3) $\alpha$鈍角，$\cos\alpha=-\dfrac{5}{6}$：$\sin\alpha=\dfrac{\sqrt{11}}{6}$，$\tan\alpha=-\dfrac{\sqrt{11}}{5}$
}

\mondai{263}{$\tan\alpha$から$\sin\alpha$，$\cos\alpha$を求めよ．}
\kotae{
$1+\tan^{2}\alpha=\dfrac{1}{\cos^{2}\alpha}$ を使う．

(1) $\alpha$鋭角，$\tan\alpha=\dfrac{1}{3}$

$\cos^{2}\alpha=\dfrac{9}{10}$，$\cos\alpha=\dfrac{3\sqrt{10}}{10}$，$\sin\alpha=\dfrac{\sqrt{10}}{10}$

(2) $\alpha$鈍角，$\tan\alpha=-2$

$\cos^{2}\alpha=\dfrac{1}{5}$，$\cos\alpha=-\dfrac{\sqrt{5}}{5}$，$\sin\alpha=\dfrac{2\sqrt{5}}{5}$
}

\mondai{264}{$\triangle$ABCにおいて，正弦定理を用いて求めよ．}
\kotae{
$\dfrac{a}{\sin A}=\dfrac{b}{\sin B}=\dfrac{c}{\sin C}=2R$

(1) $b=4$，$A=30°$，$B=45°$

$\dfrac{a}{\sin30°}=\dfrac{4}{\sin45°}$，$a=\dfrac{4\sin30°}{\sin45°}=\dfrac{2}{\frac{\sqrt{2}}{2}}=2\sqrt{2}$

(2) $b=2$，$c=\sqrt{3}$，$B=45°$

$\dfrac{2}{\sin45°}=\dfrac{\sqrt{3}}{\sin C}$，$\sin C=\dfrac{\sqrt{3}\cdot\frac{\sqrt{2}}{2}}{2}=\dfrac{\sqrt{6}}{4}$

(3) $c=5$，$A=45°$，$B=105°$

$C=30°$．$\dfrac{a}{\sin45°}=\dfrac{5}{\sin30°}=10$，$a=10\cdot\dfrac{\sqrt{2}}{2}=5\sqrt{2}$
}

\mondai{265}{1辺が$a$の正三角形の外接円の半径を求めよ．}
\kotae{
$2R=\dfrac{a}{\sin60°}=\dfrac{2a}{\sqrt{3}}$ → $R=\dfrac{\sqrt{3}\,a}{3}$
}

\mondai{266}{余弦定理を用いて辺の長さを求めよ．}
\kotae{
$a^{2}=b^{2}+c^{2}-2bc\cos A$

(1) $b=4$，$c=\sqrt{3}$，$A=30°$

$a^{2}=16+3-8\sqrt{3}\cdot\dfrac{\sqrt{3}}{2}=19-12=7$，$a=\sqrt{7}$

(2) $a=\sqrt{3}$，$c=\sqrt{6}$，$B=135°$

$b^{2}=3+6-2\sqrt{18}\cdot\left(-\dfrac{\sqrt{2}}{2}\right)=9+6=15$，$b=\sqrt{15}$

(3) $a=3\sqrt{3}$，$b=2$，$C=150°$

$c^{2}=27+4-12\sqrt{3}\cdot\left(-\dfrac{\sqrt{3}}{2}\right)=31+18=49$，$c=7$
}

\mondai{267}{$a=2$，$b=4$，$c=5$のとき$\cos A$，$\cos B$，$\cos C$を求めよ．}
\kotae{
$\cos A=\dfrac{b^{2}+c^{2}-a^{2}}{2bc}=\dfrac{16+25-4}{40}=\dfrac{37}{40}$

$\cos B=\dfrac{a^{2}+c^{2}-b^{2}}{2ac}=\dfrac{4+25-16}{20}=\dfrac{13}{20}$

$\cos C=\dfrac{a^{2}+b^{2}-c^{2}}{2ab}=\dfrac{4+16-25}{16}=-\dfrac{5}{16}$
}

\mondai{268}{$\triangle$ABCの面積を求めよ．}
\kotae{
$S=\dfrac{1}{2}bc\sin A$

(1) $b=5$，$c=7$，$A=45°$：$S=\dfrac{35}{2}\cdot\dfrac{\sqrt{2}}{2}=\dfrac{35\sqrt{2}}{4}$

(2) $a=2$，$b=3$，$C=150°$：$S=\dfrac{1}{2}\cdot2\cdot3\cdot\dfrac{1}{2}=\dfrac{3}{2}$
}

\mondai{269}{$a=7$，$B=30°$，面積9のとき$c$を求めよ．}
\kotae{
$9=\dfrac{1}{2}\cdot7\cdot c\cdot\sin30°=\dfrac{7c}{4}$ → $c=\dfrac{36}{7}$
}

\mondai{270}{$a=5$，$b=6$，$c=9$のとき．}
\kotae{
(1) $\cos C=\dfrac{25+36-81}{60}=-\dfrac{1}{3}$

(2) $\sin C=\sqrt{1-\dfrac{1}{9}}=\dfrac{2\sqrt{2}}{3}$

(3) $S=\dfrac{1}{2}\cdot5\cdot6\cdot\dfrac{2\sqrt{2}}{3}=10\sqrt{2}$
}

\mondai{271}{ヘロンの公式で面積を求めよ．}
\kotae{
$S=\sqrt{s(s-a)(s-b)(s-c)}$，$s=\dfrac{a+b+c}{2}$

(1) $a=5$，$b=7$，$c=8$：$s=10$

$S=\sqrt{10\cdot5\cdot3\cdot2}=\sqrt{300}=10\sqrt{3}$

(2) $a=2$，$b=3$，$c=4$：$s=\dfrac{9}{2}$

$S=\sqrt{\dfrac{9}{2}\cdot\dfrac{5}{2}\cdot\dfrac{3}{2}\cdot\dfrac{1}{2}}=\dfrac{3\sqrt{15}}{4}$
}

\end{document}
